\section{Convolution}

\subsection{Convolution Integral}

\subsubsection{Continuous-time convolution}
\[
(f*g)(t)=\int_{-\infty}^{\infty}f(\tau)g(t-\tau)\,d\tau.
\]
Convolution combines an input with an impulse response to produce the zero-state output of an LTIC system.

\subsubsection{Causal convolution}
\[
(f*g)(t)=\int_{0}^{t}f(\tau)g(t-\tau)\,d\tau
\quad \text{for causal }f,g.
\]
When both signals are causal, the integration interval collapses to the overlap between \(0\) and \(t\).

\subsubsection{Convolution properties}
\[
\begin{aligned}
f*g &= g*f,\\
f*(g*h)&=(f*g)*h,\\
f*(g+h)&=f*g+f*h,\\
f*\delta&=f,\qquad f*\delta(t-t_0)=f(t-t_0).
\end{aligned}
\]
Convolution is commutative, associative, and distributive. The impulse is the identity element for convolution.

\subsubsection{Differentiation and integration through convolution}
\[
\frac{d}{dt}(f*g)=\frac{df}{dt}*g=f*\frac{dg}{dt}, \qquad
u*f=\int_{-\infty}^{t}f(\tau)\,d\tau.
\]
Differentiation can be applied to either convolution factor. Convolution with the unit step integrates the signal.

\subsection{Reusable Convolution Pairs}

\subsubsection{Exponentials with equal exponents}
\[
\left(e^{-at}u(t)\right)*\left(e^{-at}u(t)\right)
=t e^{-at}u(t).
\]
Repeated first-order factors in the Laplace domain produce polynomial factors in time.

\subsubsection{Exponentials with different exponents}
\[
\left(e^{-at}u(t)\right)*\left(e^{-bt}u(t)\right)
=\frac{e^{-bt}-e^{-at}}{a-b}u(t), \qquad a\neq b.
\]
This pair is a compact exam lookup for cascaded first-order causal modes.

\subsubsection{Causal exponential and step}
\[
\left(e^{-at}u(t)\right)*u(t)=\frac{1-e^{-at}}{a}u(t), \qquad a\neq 0.
\]
This formula converts a first-order impulse response into a step response.

\subsubsection{Causal exponential and ramp}
\[
\left(e^{-at}u(t)\right)*r(t)
=\left(\frac{t}{a}-\frac{1}{a^2}+\frac{e^{-at}}{a^2}\right)u(t).
\]
This formula converts a first-order impulse response into a ramp response.

\subsubsection{Graphical convolution support}
\[
\operatorname{supp}(f*g)=\operatorname{supp}(f)+\operatorname{supp}(g).
\]
The output is nonzero only where the shifted signals overlap. For finite pulses, the breakpoints are sums of the original interval endpoints.
